\section{Related Literature}\label{sec:literature}

\subsection{Competition for productive public inputs}

Fiscal competition can distort both the level and composition of public expenditure. \citet{KeenMarchand1997} show a bias toward public inputs benefiting mobile production and away from goods valued by immobile residents, while \citet{BayindirUpmann1998} demonstrate that interjurisdictional competition over industrial public goods can generate overprovision. Firm mobility does not impose a universal sign on this distortion: \citet{Matsumoto1998} identifies underprovision under a different tax-competition environment, a distinction discussed by \citet{Ogawa2010}. Related provisions and cross-border effects also depend on market integration and mobility \citep{Ihara2008}. Excessive productive public spending, therefore, is not itself a new prediction of our analysis.

More closely, \citet{WalzWellisch1996} and \citet{MaurerWalz2000} study local public inputs and the endogenous location of oligopolistic firms. \citet{Billon2023} derives how jurisdictional market power affects tax competition when governments provide productive public inputs; \citet{Pauser2026} examines efficient provision and private utilization of congestible public inputs under unemployment and tax competition. \citet{MoritaOkoshi2025} shows that competition for foreign firms through transport infrastructure can have welfare effects even without a change in the foreign firm's location. These models establish important links between mobile firms, privately chosen activities, and public-input provision. They do not, in the cases studied, jointly determine risk-reducing public protection and privately financed geographic contingency readiness, or derive the change in a decentralized--coordinated protection gap when both backup and primary-location responses operate. That narrower comparative static, rather than a general overprovision claim, is the relevant distinction.

\subsection{Public reliability and private resilience}

Publicly supplied reliability and private substitutes have already been analyzed directly. \citet{BrehmJohnstonMilton2024} document household adoption of backup electricity technologies after outages and model the efficient provision of grid reliability in the presence of batteries and generators. Private substitutes reduce the efficient level of grid-reliability spending and affect the distribution of welfare. Accordingly, neither substitution between public reliability and private backup nor its importance for public investment appraisal is new here. Their setting concerns centralized grid-reliability provision and household adoption rather than competition between jurisdictions to attract mobile producers. Our distinct question is how protection-induced changes in firms' readiness and primary locations affect the equilibrium gap between competing governments and a policy-coordinating planner.

A separate literature studies the private--social value of production continuity. \citet{GrossmanHelpmanLhuillier2023} identify surplus and market-power wedges in firms' geographic supply-chain diversification. \citet{Eisenack2014} considers private adaptation when market structure is endogenous, and \citet{GnutzmannEtAl2020} relates spare production capacity to resilience. Such work motivates allowing firms' readiness choices to respond to policy and recognizing that their private payoffs need not capture the industrial value of availability. It does not imply that every decline in private readiness reduces welfare: lower readiness also saves real costs. In our model, the sign is determined only after both the lost availability surplus and the saved private expenses are accounted for.

Public disaster protection can also change location and exposure. \citet{KouskyLuttmerZeckhauser2006} analyze government protection and private investment in risky locations, while \citet{SeshimoYamazaki2018} show how anticipated disaster relief can encourage residential exposure and excessive prevention. Public programs can alter private defense and insurance \citep{MahmudBarbier2016}; protective infrastructure can affect firms' defense and siting \citep{GramesEtAl2019}; and the safe-development literature examines induced exposure and preparedness \citep{HaerEtAl2020,FusinatoEtAl2024}. \citet{GawelEtAl2018} discusses public-choice obstacles to efficient adaptation. These studies are relevant precedents for behavioral responses to protection, but they do not establish our particular comparison of globally optimal unilateral protection choices with re-solved private production continuity and location decisions.

Public--private protective substitution also arises outside disaster economics. \citet{HickeyEtAl2021} distinguish technological from equilibrium responses in private protection and public policing; \citet{FrieheMungan2022} studies political overprovision of public policing; and \citet{ChoiBaeJun2010} analyzes substitution between public and private antipiracy protection. \citet{LeePinto2009} examines prevention across jurisdictions. These precedents reinforce that the novelty cannot rest on substitution or excessive public protection alone.

\subsection{Endogenous location, market structure, and the policy gap}

The interaction of industrial location and public policy has a substantial literature. \citet{MarkusenMoreyOlewiler1993,MarkusenMoreyOlewiler1995} endogenize plant location in environmental-policy competition; \citet{MartinHerranEtAl2026} studies environmental policy and industrial allocation in a dynamic spatial setting. Research on disaster exposure and production networks includes \citet{OtazawaEtAl2016}, \citet{BacchiocchiEtAl2024}, and \citet{MengDingWang2024}. \citet{HirteEtAl2018} also shows how local incentives can favor excessive adaptation relative to welfare. Their instruments, private margins, and welfare criteria differ, but together they make clear that neither strategic location nor resilience-oriented policy competition is a novel ingredient.

Against these predecessors, our contribution is deliberately specific. A first result shows that risk-reducing public inputs can lower national industrial surplus through firms' privately chosen loss of contingency readiness while remaining attractive to individual jurisdictions seeking mobile production. The stronger distinction concerns a common direct protection benefit: in an exactly solved policy game, its effect on the decentralized-minus-coordinated protection gap is positive when backup and location are endogenous, zero if backup is held fixed, and negative if location is held fixed. Each of the three policy games is fully re-solved, with the restrictions in the latter two maintained under every unilateral government deviation. General sufficient conditions and a local differentiated-product extension connect the comparison to a wider, though not universal, set of primitives. These are within-model statements about an industrial-surplus benchmark with decentralized private choices, not claims that public--private substitution, firm attraction, or public overprovision are new in general.
